\documentclass[12pt]{article}
\usepackage[utf8]{inputenc}
\usepackage[T1]{fontenc}
\usepackage[margin=1in]{geometry}
\usepackage{setspace}
\usepackage{xcolor}
\usepackage[colorlinks=true, urlcolor=blue!60!black, linkcolor=blue!60!black]{hyperref}
\setlength{\parindent}{0em}
\setlength{\parskip}{0.6em}
\renewcommand{\baselinestretch}{1.0}
\begin{document}
\begin{center}
{\LARGE \textit{Carson v. Makin}, 596 U.S. 767 (2022)\par}
\vspace{0.5em}
{\large Case Brief --- Two-Page Edition\par}
\end{center}

\section{Memory Jogger}
Maine cannot exclude ``sectarian'' schools from a generally available tuition program; religious ``use'' discrimination is no safer than ``status'' discrimination (\href{https://www.courtlistener.com/opinion/6479298/carson-v-makin/}{CourtListener}).

\section{Detailed Case Facts}
Most Maine school units run no public high school, so they pay tuition at the public or approved private school a parent selects. Approved private schools may be selective, single-sex, and free of state curricular mandates. Since 1981, Maine had limited the program to ``nonsectarian'' schools, defined as those promoting a faith or teaching through its lens. The Carson and Nelson families wanted to use the benefit at accredited Christian schools.

\section{Procedural History}
The District of Maine ruled for Maine. After \textit{Espinoza} (2020), the First Circuit still affirmed, reasoning that Maine barred religious ``use,'' not ``status,'' and funded the ``rough equivalent'' of public education. The Supreme Court reversed.

\section{Judicial Votes}
6--3. Chief Justice Roberts wrote for Thomas, Alito, Gorsuch, Kavanaugh, and Barrett. Breyer dissented, joined by Kagan and in part by Sotomayor; Sotomayor also dissented separately.

\section{Holding}
The ``nonsectarian'' requirement violates the Free Exercise Clause. Once a State funds private schooling, excluding religious schools for their status or their religious use triggers strict scrutiny, and an interest in stricter separation than the Establishment Clause requires is not compelling.

\section{Analysis of Opinions}
\textbf{Majority:} Under \textit{Trinity Lutheran} and \textit{Espinoza}, benefit exclusions penalize free exercise. Parents' independent choice satisfies \textit{Zelman}. The program did not mirror public schools, and Maine could instead expand public options. Policing religious use would breed entanglement, and \textit{Locke v. Davey} is limited to clergy training. \textbf{Breyer:} the Court ignored the Establishment Clause and erased the ``play in the joints.'' \textbf{Sotomayor:} taxpayers must now fund schools that discriminate on religion, sexual orientation, and gender identity.

\section{Examples: Future Applications}
\textbf{Same-Side:} an education savings account barring schools that integrate religion into academics is unconstitutional. \textbf{Opposite-Side:} a State may end private tuition entirely and build public schools instead. \textbf{Fence-Sitter:} a neutral requirement that all participating schools follow nondiscrimination laws in admissions and hiring, which footnote 1* left open.

\section{Critique}
Ira C. Lupu and Robert W. Tuttle argue the Court discarded the status/use line when \textit{Zelman} offered a narrower ground, unsettling decades of Establishment Clause precedent. Defenders such as Michael Bindas answer that the Clause protects religious exercise, not mere identity, and that the old line forced theological auditing.

\section{Key Quotations}
\begin{quote}
``In short, the prohibition on status-based discrimination under the Free Exercise Clause is not a permission to engage in use-based discrimination.''
\end{quote}
\noindent --- Chief Justice Roberts, for the Court, 596 U.S. 767.

\end{document}
